% ------------------------------------------------------------
% Vs'shtak-/Draconic-FDI-Schrift
% ------------------------------------------------------------
% Wird von \FDIInputLanguageFonts in der Praeambel geladen.
% Wichtig: Diese Fonts werden NICHT global auf das Dokument angewendet,
% sondern nur ueber die FDI-Font-Makros verwendet.
%
% Iokharic ist das drakonische Alphabet. Es ersetzt hier ausschliesslich
% die Overlay-Schrift, also Sprechblasen, Kaesten und Bildunterschriften.
%
% Bewusst NICHT ersetzt werden:
%   \KarolingischeMinuskel     Seite 2, Panel 5 und 6 zeigen die historische
%                              karolingische Minuskel selbst. Sie ist dort
%                              Inhalt, nicht Typografie, und bleibt unveraendert.
%                              Das ist in dieser Ausgabe die einzige Stelle.
%                              Die Karten auf Seite 1 und 3 benutzen dagegen
%                              direkt \FDIOverlayFont, dort ist die Minuskel
%                              reine Typografie. Siehe pages_vss/page_1.tex
%                              und pages_vss/page_3.tex.
%   \ArialFont, \ArialBOLDFont URLs, DOI, ISBN und Diagrammbeschriftungen
%                              muessen lesbar und maschinenlesbar bleiben.
%
% Die Schriftdateien liegen NICHT im Repository, siehe fonts_vss/FONT_NOTE.txt.
% Fehlen sie, faellt diese Datei mit einer Warnung auf die Standardschrift
% zurueck, damit das Dokument trotzdem uebersetzt.
%
% ------------------------------------------------------------
% Warum hier ueberhaupt Lua steht
% ------------------------------------------------------------
% Iokharic besitzt echte Glyphen nur fuer a-z, A-Z und 0-9. Saemtliche
% Satzzeichen, Anfuehrungszeichen, Apostrophe und Umlaute sind im Font auf
% eine identische Platzhalter-Box gemappt. LaTeX meldet das nicht als
% "missing character", weil die Glyphen formal vorhanden sind; im PDF stehen
% dann leere Kaesten. Deshalb zwei Eingriffe:
%
%   1. luaotfload.patch_font entfernt die Platzhalter aus der geladenen
%      Fonttabelle. Damit gelten sie als fehlend und der Fallback-Font
%      (die Standard-Comicschrift Komika) liefert sie nach. Die Fontdatei
%      selbst wird dabei nicht angefasst.
%   2. Der Bindestrich ist ein Sonderfall: LuaTeX baut daraus ein
%      Discretionary und verwirft das Zeichen bereits, bevor der Fallback
%      laufen kann. Er bleibt deshalb in der Tabelle stehen und wird
%      stattdessen im Knotenbaum auf den Fallback-Font umgehaengt.
%
% Beide Callbacks pruefen den Dateinamen und wirken daher nur auf Iokharic.
% Die anderen Sprachausgaben sind unberuehrt.
%
% Innerhalb von \directlua stehen bewusst keine Leerzeilen: dieser Block ist
% ein Makroargument von \IfFileExists, und eine Leerzeile darin wuerde als
% \par im Lua-Code landen.

\IfFileExists{language_files/vss/fonts_vss/Iokharic-Regular.otf}{%
\directlua{
  luaotfload.add_fallback("fdivssfallback", { "[fonts/KOMTXT__.ttf]:mode=node;" })
  FDIVss = { hyphenfont = {} }
  % 1. Platzhalterglyphen aus der geladenen Fonttabelle werfen.
  luatexbase.add_to_callback("luaotfload.patch_font", function(fontdata)
    if not (fontdata.filename or ""):find("Iokharic", 1, true) then return end
    local chars = fontdata.characters
    if not chars then return end
    local keep = {[0x20] = true, [0x2D] = true}
    for c = 0x30, 0x39 do keep[c] = true end
    for c = 0x41, 0x5A do keep[c] = true end
    for c = 0x61, 0x7A do keep[c] = true end
    for uni in next, chars do
      if not keep[uni] then chars[uni] = nil end
    end
  end, "fdi.vss.strip")
  % 2. Bindestriche aus Iokharic auf den Fallback-Font umhaengen.
  local glyph_id = node.id("glyph")
  local disc_id  = node.id("disc")
  local hlist_id = node.id("hlist")
  local vlist_id = node.id("vlist")
  % Kein "~=" verwenden: ~ ist in LaTeX aktiv und landet als Backslash im Lua.
  local function hyphenfont(fid)
    local cached = FDIVss.hyphenfont[fid]
    if cached == nil then
      local f = font.getfont(fid)
      cached = false
      if f and (f.filename or ""):find("Iokharic", 1, true) then
        local lookup = f.fallback_lookup
        cached = lookup and lookup[0x2D] or false
      end
      FDIVss.hyphenfont[fid] = cached
    end
    return cached
  end
  local function walk(head)
    local n = head
    while n do
      local id = n.id
      if id == glyph_id then
        if n.char == 0x2D then
          local repl = hyphenfont(n.font)
          if repl then n.font = repl end
        end
      elseif id == disc_id then
        walk(n.pre) walk(n.post) walk(n.replace)
      elseif id == hlist_id or id == vlist_id then
        walk(n.head)
      end
      n = n.next
    end
  end
  local function fdivssfilter(head)
    walk(head)
    return true
  end
  luatexbase.add_to_callback("pre_linebreak_filter", fdivssfilter, "fdi.vss.hyphen")
  luatexbase.add_to_callback("hpack_filter",        fdivssfilter, "fdi.vss.hyphen")
}
\newfontfamily{\FDIVssOverlayFont}{Iokharic-Regular.otf}[
  Path           = language_files/vss/fonts_vss/,
  BoldFont       = Iokharic-Bold.otf,
  ItalicFont     = Iokharic-Italic.otf,
  BoldItalicFont = Iokharic-BoldItalic.otf,
  Scale          = 1.00,
  Ligatures      = NoCommon,
  RawFeature     = {fallback=fdivssfallback}
]
% Zuweisung der FDI-Font-Makros fuer diese Sprache.
% Wird von \FDISetupLanguageFonts ausgefuehrt, sobald vss gesetzt ist.
\FDISetLanguageFonts{vss}{%
  \let\FDIOverlayFont\FDIVssOverlayFont
}%
}{%
  \PackageWarningNoLine{FDI}{%
    Iokharic not found, the vss edition falls back to the default overlay
    font.\MessageBreak
    See language_files/vss/fonts_vss/FONT_NOTE.txt%
  }%
}
